\Needspace{11\baselineskip}
\subsection{A positive-volume endpoint that is not big}
\label{subsec:actual-flow-positive-volume-nonbig}

Backgrounds reached at the maximal time of a Chern--Ricci flow are not better behaved, as far as the relation between positive volume and bigness is concerned. The radial background of \cref{sec:radial-positive-volume-nonbig} occurs, up to a smooth $\ddc$ change, as the limiting background of a Chern--Ricci flow at its maximal time. The normalized currents annihilating it can also be described explicitly.

We keep the notation of \cref{sec:radial-positive-volume-nonbig}: $X=\mathbb P^m\times\mathbb P^1$ with $n=m+1\geq3$, the normalized Fubini--Study forms $\omega$ and $w$ pulled back from the two factors, the nonnegative function $a=x/t$ on the first factor, which vanishes only at $p$, the background $\theta=\omega+aw$ and the probability form $\mu=\omega^m\wedge w$. Put $Z=\{p\}\times\mathbb P^1$, and start the flow from the Hermitian metric
\begin{equation}\label{eq:radial-flow-initial-metric}
 \omega_0=(n+1)\omega+(a+2)w.
\end{equation}

\begin{theorem}\label{thm:actual-flow-positive-volume-nonbig}\index{Chern--Ricci flow!nonbig positive-volume endpoint}
With the normalization of the Ricci form in \cref{subsec:chern-ricci-maximal-time}, the maximal smooth existence time of the Chern--Ricci flow starting from $\omega_0$ is $\tau=1/(2\pi)$. The limiting background $\alpha_\tau=\omega_0-\tau\operatorname{Ric}(\omega_0)$ is nef and not big, and
\[
 \alpha_\tau=\theta+\ddc\log(a+2),\qquad
 \mathcal V_\mu(\alpha_\tau)=n.
\]
Its variational volume is positive for every smooth strictly positive probability form on $X$.
\end{theorem}
\begin{proof}
The initial metric is strictly positive, and
\[
 \omega_0^n=n(n+1)^m(a+2)\mu.
\]
With our normalization, the Ricci forms of the normalized Fubini--Study metrics are $2\pi n\omega$ and $4\pi w$. Consequently
\[
 \operatorname{Ric}(\omega_0)
   =2\pi\bigl(n\omega+2w-\ddc\log(a+2)\bigr).
\]
Write $s$ for the flow time and put $r=2\pi s$. The affine background is
\begin{equation}\label{eq:radial-flow-affine-background}
 \alpha_s=(n+1-nr)\omega+(a+2-2r)w
                          +r\ddc\log(a+2).
\end{equation}
For $0\leq r<1$, adding $\ddc$ of the smooth function $-r\log(a+2)$ gives a strictly positive representative: the horizontal coefficient is at least one and the vertical coefficient is at least $2(1-r)>0$. At $r=1$ we obtain the identity $\alpha_\tau=\theta+\ddc\log(a+2)$. Nefness, bigness and the variational volume are unchanged by this smooth $\ddc$ change, by the conventions of \cref{chap:gauduchon-variational-volume} and \cref{prop:gaud-volume-elementary}, so \cref{thm:radial-positive-volume-nonbig} shows that the endpoint background is nef, not big, and of volume $n$.

For $r>1$, choose a nonzero smooth nonnegative function $f$ on the base supported in $\{a+2-2r<0\}$, a neighborhood of $p$. Its pullback $S=f\omega^m$ is a closed semipositive $(n-1,n-1)$-form. For every quasi-plurisubharmonic function $v$, Stokes' theorem for currents gives
\[
 \int_X(\alpha_s+\ddc v)\wedge S
       =\int_{\mathbb P^m}f(a+2-2r)\omega^m<0,
\]
so $\alpha_s$ is not even pseudoeffective for $r>1$. At $r=1$, a smooth positive representative would contradict nonbigness. Thus $r=1$ is the supremum of the admissible times, and \cref{thm:chern-ricci-maximal-time} gives the smooth flow on $[0,\tau)$ with exactly this lifetime; the reference representatives used above need not solve the flow equation. Finally, \cref{prop:gaud-volume-density-comparison} gives positivity of the volume for every smooth positive probability form.
\end{proof}

\begin{proposition}\label{prop:radial-flow-annihilator-support}\index{Gauduchon current!flow endpoint support}
The current $[Z]/2$ belongs to $\mathscr G_{\omega_0}$ and annihilates $\alpha_\tau$. Every $R\in\mathscr G_{\omega_0}$ with $\int_X\alpha_\tau\wedge R=0$ has its trace probability $\sigma_R=\omega_0\wedge R$ supported on $Z$. In the coefficient-matrix convention of \cref{sec:weak-gauduchon-duality}, its matrix is
\[
 P_R=n\Pi_{\mathrm{vert}}\qquad\sigma_R\text{-almost everywhere},
\]
where $\Pi_{\mathrm{vert}}$ is the $\omega_0$-orthogonal projection onto the tangent line of the fiber $\{p\}\times\mathbb P^1$. In particular, $P_R$ has rank one, and the determinant measure of $R$ vanishes for every smooth positive density.
\end{proposition}
\begin{proof}
Only smooth forms need to be restricted to $Z$, and their restrictions satisfy $\omega_0|_Z=2w$ and $\alpha_\tau|_Z=0$. This proves that $[Z]/2$ is normalized and annihilates $\alpha_\tau$.

Let $R$ be a general annihilator. Since $\alpha_\tau$ and $\theta$ differ by $\ddc$ of a smooth function, Stokes' theorem gives $\int_X\theta\wedge R=0$. The measure $\theta\wedge R$ is nonnegative, so it vanishes. On every compact set $K\subset X\setminus Z$, the explicit form $\theta$ satisfies $\theta\geq c_K\omega_0$ for some $c_K>0$. Thus $\sigma_R(K)=0$, and a countable exhaustion by compact sets shows that $\sigma_R$ is supported on $Z$.

Relative to $\omega_0$, the matrix of $R$ is characterized, for smooth test forms $\beta$ with matrix $B_\beta$, by
\[
 \beta\wedge R=\frac1n\operatorname{tr}(B_\beta P_R)\sigma_R,
 \qquad P_R\geq0,\qquad\operatorname{tr}P_R=n.
\]
At points of $Z$, the form $\theta$ is positive on the horizontal $m$-plane and vanishes on the vertical line. The vanishing of $\theta\wedge R$ therefore forces all horizontal diagonal entries of $P_R$ to vanish, and the inequality $|(P_R)_{ij}|^2\leq(P_R)_{ii}(P_R)_{jj}$ for positive matrices removes their rows and columns. The remaining vertical entry is $n$, which proves the formula and the vanishing of the determinant.
\end{proof}

For each such annihilator $R$, the trace identity at the end of \cref{subsec:chern-ricci-maximal-time} becomes
\[
 \int_X\omega(s)\wedge R=1-2\pi s,
 \qquad 0\leq s<1/(2\pi).
\]
The support statement uses the explicit semipositive representative $\theta$ of the endpoint class; the vanishing of the determinant alone would not give it.
